\subsection*{1.4}\phantomsection \label{sga3:VIIB:1.4}\addcontentsline{toc}{subsection}{1.4. 下降与正合性}

\noindent\textbf{定理 1.4。}\phantomsection \label{sga3:VIIB:1.4.thm} {\itshape 设 \(k\) 为伪紧环，\[
  d_0,d_1:X_1\rightrightarrows X
\] 为 \(\operatorname{Vaf}/k\) 中的等价对 \textup{（参见第五讲，\S~2.b）}，且 \(d_1\) 拓扑平坦。}\begin{enumerate}\item[\textup{(i)}]{\itshape 从 \(X\) 到 \(X/X_1=\operatorname{Coker}(d_0,d_1)\) 的典范投影满射且拓扑平坦，以 \(d_0\) 和 \(d_1\) 为分量的态射 \[
    X_1\longrightarrow X\times_{X/X_1}X
  \] 是同构。}

\item[\textup{(ii)}]{\itshape 若 \(X\) 在 \(k\) 上拓扑平坦，则 \(X/X_1\) 也如此。}\end{enumerate}

首先注意，由 \SGARefLink{sga3:VIIB:1.3.3}{1.3.3}\,\textup{\SGARefLink{sga3:vii:paragraph:5-1:equation:ii}{(ii)}} 可知 \textup{\SGARefLink{sga3:vii:paragraph:5-1:equation:i}{(i)}} 蕴含 \textup{\SGARefLink{sga3:vii:paragraph:5-1:equation:ii}{(ii)}}。\textup{\SGARefLink{sga3:vii:paragraph:5-1:equation:i}{(i)}} 的证明见 \SGARefLink{sga3:VIIA:1.4.1}{1.4.1}、\SGARefLink{sga3:VIIA:1.4.2}{1.4.2} 和 \SGARefLink{sga3:VIIB:1.4.3}{1.4.3}。

\subsubsection*{1.4.1}
\phantomsection
\label{sga3:VIIB:1.4.1}

先说明可以归约到 \(X\) 只有一个点的情形。由于这里处理的是等价对，与第五讲 \S~3.e) 一样，可以在 \(X\) 的底集合上定义等价关系：当存在点 \(z\in X_1\) 使得 \[
  d_0(z)=y
  \qquad\text{且}\qquad
  d_1(z)=x.
\] 时，称两点 \(x,y\) 等价。显然不失一般性，可以假设 \(X\) 对此关系只有一个等价类；换言之，\(X/X_1\) 只有一个点 \textup{（见 \SGARefLink{sga3:VIIB:1.2}{1.2} 中给出的 \(X/X_1\) 的构造）}。

此时，设 \(x\) 为 \(X\) 的任一点，\(U\) 为以 \(x\) 为唯一点、且局部环与 \(X\) 在 \(x\) 处的局部环相同的形式簇。与第五讲 \S~6 一样可知，\((d_0,d_1)\) 在 \(U\) 上诱导的等价关系仍满足定理的假设，且只须对此等价关系证明定理 \textup{（\(U\) 是一个“拟截面”）}。

简要回顾第五讲 \S~6 中证明的原理。置 \[
  V=d_0^{-1}(U)=U\times_{i,d_0}X_1,
\]，其中 \(i\) 为 \(U\) 到 \(X\) 的包含映射，设 \(u\) 和 \(v\) 分别为 \(d_0\) 和 \(d_1\) 诱导的以 \(V\) 为源的态射：\begin{center}\phantomsection \label{sga3:VIIB:1.4.1:quasisection-span} % SGA3-VII-DIAG-090
% Authority: Exp7B-13oct24.pdf, local p. 35.
% Crop: figures/Exp7B_localp035_quasisection_span.png
% Locator: tex/components/29_expose_VIIB_section14_descent_exactness_en.tex, Loop 1 line 82.
\(
X\xleftarrow{\ v\ }V\xrightarrow{\ u\ }U.
\)
\end{center}显然 \(u\) 和 \(v\) 拓扑平坦，且 \(u\) 满射；由于 \(X\) 只有一个等价类，\(v\) 也满射。若 \((v_0,v_1)\) 是等价对 \((d_0,d_1)\) 在 \(v\) 下的逆像 \textup{（参见 V，3.a)）}，则由 V，3.c) 和 3.d) 可知，\(X/X_1\) 和 \(U\) 对 \((d_0,d_1)\) 诱导的等价关系的商均可与 \(\operatorname{Coker}(v_0,v_1)\) 等同。与 V，\SGARefLink{lem:sga3-V-6.1}{6.1} 的证明一样可知，若 \SGARefLink{sga3:VIIB:1.4.thm}{定理~1.4} 的结论对 \(U\) 成立，则对 \(X\) 也成立。

\subsubsection*{1.4.2}
\phantomsection
\label{sga3:VIIB:1.4.2}

这样就归约到 \(X\) 只有一个点的情形。\footnote{\label{sga3:VIIB:nde:67}编者注：因此可以假设 \(k\) 为局部环。}考虑下面的交换图 \textup{（参见 V，\S~1，\textup{(0,1,2)}）}：\begin{center}\phantomsection \label{sga3:VIIB:1.4.2:equivalence-groupoid-diagram} % SGA3-VII-DIAG-091
% Authority: Exp7B-13oct24.pdf, local p. 35.
% Crop: figures/Exp7B_localp035_equivalence_groupoid_diagram.png
% Locator: tex/components/29_expose_VIIB_section14_descent_exactness_en.tex, Loop 1 line 107.
\begin{tikzcd}[row sep=3.3em,column sep=3.9em]
X_2
  \arrow[r,shift left=.6ex,"d'_1"]
  \arrow[r,shift right=.6ex,"d'_0"']
  \arrow[d,"d'_2"']
  & X_1 \arrow[r,"d_0"] \arrow[d,"d_1"]
  & X \arrow[d] \\
X_1
  \arrow[r,shift left=.6ex,"d_1"]
  \arrow[r,shift right=.6ex,"d_0"']
  & X \arrow[r]
  & X/X_1
\end{tikzcd}
\end{center}这里 \[
  X_2=X_1\times_{d_1,d_0}X_1,
\]，而 \(d'_0,d'_1,d'_2\) 分别为态射\footnote{\label{sga3:VIIB:nde:68}编者注：参见第五讲，\S~2.b)。}\[
  (x,y,z)\longmapsto(x,y),\qquad
  (x,y,z)\longmapsto(x,z),\qquad
  (x,y,z)\longmapsto(y,z).
\]。

若以 \(B,A,A_1,A_2\) 分别表示 \(X/X_1,X,X_1,X_2\) 的仿射代数，则上述图诱导交换图\begin{center}\phantomsection \label{sga3:VIIB:1.4.2:affine-algebra-diagram} % SGA3-VII-DIAG-092
% Authority: Exp7B-13oct24.pdf, local p. 35.
% Crop: figures/Exp7B_localp035_affine_algebra_diagram.png
% Locator: tex/components/29_expose_VIIB_section14_descent_exactness_en.tex, Loop 1 line 129.
\begin{tikzcd}[row sep=3.3em,column sep=4.1em]
A_2
  & A_1
      \arrow[l,shift left=.6ex,"j_0"]
      \arrow[l,shift right=.6ex,"j_1"']
  & A \arrow[l,"i_0"'] \\
A_1 \arrow[u,"j_2"]
  & A
      \arrow[l,shift left=.6ex,"i_0"]
      \arrow[l,shift right=.6ex,"i_1"']
      \arrow[u,"i_1"']
  & B \arrow[l,"i"'] \arrow[u,"i"]
\end{tikzcd}
\end{center}，其中两行均正合，由 \((i_0,j_0)\) 和 \((i_1,j_1)\) 确定的方块均余笛卡尔。由假设，以 \(d_0\) 和 \(d_1\) 为分量的态射 \[
  X_1\longrightarrow X\times X
\] 是单态射，故由 \SGARefLink{sga3:VIIB:1.3.prop}{命题~1.3}，以 \(i_0\) 和 \(i_1\) 为分量的态射 \[
  A\widehat{\otimes}_k A\longrightarrow A_1
\] 满射。

\begin{samepage}这意味着 \(i_1\) 使 \(A_1\) 成为由 \(i_0(A)\) 生成的伪紧 \(A\)-模，且由假设它拓扑自由。由于 \(A\) 为局部环，下面的 \SGARefLink{sga3:VIIB:1.4.3}{引理~1.4.3} 给出拓扑自由 \(k\)-模 \(V\) 以及伪紧 \(k\)-模态射 \(f:V\to A\)，使得态射 \[
  \alpha_1:
  A\widehat{\otimes}_kV\longrightarrow A_1,
  \qquad
  a\widehat{\otimes}v
  \longmapsto
  i_1(a)\cdot i_0\bigl(f(v)\bigr)
\] 为同构。\end{samepage}\begin{samepage}设 \[
  \alpha:B\widehat{\otimes}_kV\longrightarrow A
  \qquad\text{且}\qquad
  \alpha_2:A_1\widehat{\otimes}_kV\longrightarrow A_2
\] 为态射 \[
  b\widehat{\otimes}v
  \longmapsto
  i(b)\cdot f(v)
  \qquad\text{且}\qquad
  a_1\widehat{\otimes}v
  \longmapsto
  j_2(a_1)\cdot j_0i_0\bigl(f(v)\bigr).
\]。\end{samepage}在下面的交换图中，两行均正合，左侧两个方块均余笛卡尔。由于 \(\alpha_1\) 为同构，\(\alpha_2\) 亦然，进而 \(\alpha\) 也是同构。\begin{center}\phantomsection \label{sga3:VIIB:1.4.2:alpha-cocartesian-diagram} % SGA3-VII-DIAG-093
% Authority: Exp7B-13oct24.pdf, local p. 36.
% Crop: figures/Exp7B_localp036_alpha_cocartesian_diagram.png
% Locator: tex/components/29_expose_VIIB_section14_descent_exactness_en.tex, Loop 1 line 186.
\begin{tikzcd}[row sep=3.4em,column sep=3.5em]
A_2
  & A_1
      \arrow[l,shift left=.6ex,"j_0"]
      \arrow[l,shift right=.6ex,"j_1"']
  & A \arrow[l,"i_0"'] \\
A_1\widehat\otimes_k V \arrow[u,"\alpha_2"]
  & A\widehat\otimes_k V
      \arrow[l,shift left=.6ex,"{i_0\widehat\otimes V}"]
      \arrow[l,shift right=.6ex,"{i_1\widehat\otimes V}"']
      \arrow[u,"\alpha_1"']
  & B\widehat\otimes_k V
      \arrow[l,"{i\widehat\otimes V}"']
      \arrow[u,"\alpha"']
      \, .
\end{tikzcd}
\end{center}

这一方面说明 \(A\) 在 \(B\) 上拓扑自由，伪基为 \(f(V)\) \textup{（参见 \SGARefLink{sga3:VIIB:0.2.1}{0.2.1}）}；另一方面说明，将 \(A_1\) 通过 \(i_1\) 视为 \(A\)-模时，\(A_1\) 在 \(A\) 上的一个伪基由 \(f(V)\) 在 \(i_0\) 下的像给出。因此，以 \(i_1\) 和 \(i_0\) 为分量的态射 \[
  A\widehat{\otimes}_B A\longrightarrow A_1
\] 为同构：\[
  A\widehat{\otimes}_B A
  \simeq
  A\widehat{\otimes}_B B\widehat{\otimes}_kV
  \simeq
  A\widehat{\otimes}_kV
  \simeq
  A_1.
\]。除下面的引理外，\SGARefLink{sga3:VIIB:1.4.thm}{定理~1.4} 已证。

\subsubsection*{1.4.3}
\phantomsection
\label{sga3:VIIB:1.4.3}

\begin{samepage}\noindent\textbf{引理 1.4.3。}{\itshape 设 \(k\) 为伪紧环，\(A\) 为局部预有限 \(k\)-代数，\(A_1\) 为拓扑自由 \(A\)-模，\[
  i_0:M\longrightarrow A_1
\] 为伪紧 \(k\)-模态射。假设映射 \[
  A\widehat{\otimes}_kM\longrightarrow A_1,
  \qquad
  a\widehat{\otimes}m
  \longmapsto
  a\cdot i_0(m),
\] 满射。则存在拓扑自由 \(k\)-模 \(V\) 和伪紧 \(k\)-模态射 \[
  f:V\longrightarrow M
\]，使得映射 \[
  A\widehat{\otimes}_kV\longrightarrow A_1,
  \qquad
  a\widehat{\otimes}v
  \longmapsto
  a\cdot i_0\bigl(f(v)\bigr),
\] 为双射。}\end{samepage}

由于每个伪紧 \(k\)-模都是拓扑自由 \(k\)-模的商 \textup{（参见编者注 \textup{(27)}）}，不失一般性可以假设 \(M\) 拓扑自由。写成 \[
  M=\prod_{i\in I}M_i,
\]，其中每个 \(M_i\) 都是 \(k\) 的一个副本。此时 \[
  A\widehat{\otimes}_kM
  =
  \prod_{i\in I}
  A\widehat{\otimes}_kM_i.
\]。由于映射 \[
  a\widehat{\otimes}m\longmapsto a\cdot i_0(m)
\] 满射且 \(A_1\) 为投射模，其核是 \(A\widehat{\otimes}_kM\) 的直因子。由于 \(A\) 为局部环，交换定理 \textup{（\SGARefLink{sga3:VIIB:0.3.4}{0.3.4}）} 表明，此核有一个补模，是部分乘积 \[
  \prod_{i\in J}A\widehat{\otimes}_kM_i,
\]，其中 \(J\) 为 \(I\) 的子集。因此可以取 \[
  V=\prod_{i\in J}M_i.
\]。


